\section{El CEO en la transición: custodio de la marca, decisiones 51/49 y espíritu emprendedor}

Las dos secciones anteriores trataron el mercado y la tecnología; esta se ocupa de la organización y del CEO, porque la transformación no la ejecuta por sí sola una hoja de ruta tecnológica. La pregunta que se responde aquí es la siguiente: cuando una empresa tiene 139 años de historia, 160,000 empleados, presencia en el mercado global y una enorme base de clientes existente, ¿cómo puede el CEO mantener la ofensiva sin destruir la credibilidad de la marca? Cuando Kang Linsong (康林松) asumió como CEO en 2019, ya llevaba muchos años trabajando en Mercedes-Benz; sin embargo, al convertirse realmente en CEO, «al mirar atrás, no hay nadie más, solo tú». Esta expresión es fundamental: en un gran sistema con 160,000 empleados y presencia en más de 150 países, la responsabilidad final recae en el CEO, y las decisiones durante una transformación no suelen ser 80/20, sino 51/49.

\begin{figure}[H]
\centering
\includegraphics[width=0.96\textwidth]{transformation-ceo-loop.png}
\caption{Ciclo de decisión del CEO en la transición: decisiones 51/49, aprendizaje continuo y corrección cuando sea necesario. Diagrama conceptual propio, elaborado a partir de la conversación de 00:40:17--00:50:00.}
\end{figure}

\begin{knowledgebox}{Lectura de la figura: en una transición, decidir no es una apuesta única}
En la figura, Uncertainty, Decision, Observe, Correct y Guardian forman un ciclo cerrado. Kang Linsong subraya que, si más adelante se descubre que la elección fue errónea, hay que recopilar información nueva y corregir, sin quedar atrapado por el orgullo ni por los costos hundidos.
\end{knowledgebox}

\subsection{Cómo conservan el espíritu emprendedor los directivos profesionales}

La sección anterior mostró que la responsabilidad del CEO no puede transferirse; esta examina cómo una empresa madura evita que su organización se vuelva lenta. Zhang Xiaojun (张小珺) preguntó si, dado que Tesla y los fabricantes de automóviles chinos siguen impulsados por sus fundadores, mientras que los fabricantes alemanes han pasado por varias generaciones de directivos profesionales, estos últimos son más conservadores. La respuesta de Kang Linsong fue que una empresa madura también debe conservar el espíritu emprendedor, a la vez que enfrenta la presión de los mercados financieros y la de la transformación. No puede limitarse a jugar a lo seguro ni únicamente a defenderse. Debe invertir en nuevas tecnologías, aumentar la eficiencia y simplificar su estructura, para que la organización se mantenga en condición competitiva.

\begin{warningbox}{La doble presión de la transformación en una gran empresa}
Una empresa emergente puede cambiar pérdidas por velocidad; una empresa madura que cotiza en bolsa enfrenta al mismo tiempo las exigencias de utilidades, accionistas, marca, empleados e inversión en la transformación. La dificultad para el CEO en la transición consiste en preservar la solidez financiera sin perder la ventana tecnológica por jugar a lo seguro.
\end{warningbox}

\subsection{La tensión entre simplificación, velocidad y una estrategia de muchos modelos}

La sección anterior planteó que una empresa madura debe conservar el espíritu emprendedor; esta entra en una contradicción concreta de gestión. Kang Linsong insiste en la simplificación, la velocidad y la eficiencia, y al mismo tiempo afirma que en los próximos tres años Mercedes-Benz lanzará más modelos que nunca. Parece una contradicción, pero en realidad es una relación entre ataque y defensa: los procesos de la organización deben simplificarse, pero la ofensiva de producto no puede detenerse; los costos deben controlarse, pero los delanteros también tienen que entrar a la cancha. Según su metáfora futbolística, no se gana un partido solo con el portero, y una empresa madura no puede proteger sus utilidades únicamente defendiéndose.

\noindent\begin{tabularx}{\textwidth}{p{0.20\textwidth}p{0.34\textwidth}X}
\toprule
Objetivo de gestión & Interpretación errónea frecuente & Interpretación más precisa \\
\midrule
Simplificación & Hacer menos productos, reducir la inversión & Reducir niveles jerárquicos y desperdicio en los procesos, y acelerar las decisiones. \\
Eficiencia & Solo recortar costos & Organizar de forma más ágil la investigación y desarrollo, las plantas, la cadena de suministro y los socios. \\
Estrategia de muchos modelos & Contradice la simplificación & La ofensiva de producto debe mantenerse; la complejidad organizacional debe gestionarse. \\
Espíritu emprendedor & Ignorar las utilidades como una empresa emergente & Mantener la ofensiva y la inversión en nuevas tecnologías dentro de las restricciones financieras. \\
\bottomrule
\end{tabularx}

\subsection{El custodio de la marca}

Lo anterior trató cómo actúa la organización; esta sección vuelve a cómo define el éxito el CEO. Para Kang Linsong, una transformación exitosa no se define por una sola cifra de participación de mercado, de margen de utilidad o de liderazgo tecnológico, sino por actuar como custodio de la marca Mercedes-Benz y entregarla a su sucesor con la marca, la experiencia del cliente, la solidez financiera y la capacidad de innovación en mejor estado. Esta definición explica bien la escala temporal de una empresa industrial centenaria: la tecnología cambia, los motores cambian, la energía cambia, pero la promesa de la marca debe perdurar.

\begin{figure}[H]
\centering
\includegraphics[width=0.94\textwidth]{guardian-of-brand.png}
\caption{El custodio de la marca: el éxito no es un indicador técnico aislado, sino que la marca sea más fuerte al momento del relevo. Diagrama conceptual propio, elaborado a partir de la conversación de 00:51:03--00:57:57.}
\end{figure}

\begin{knowledgebox}{Lectura de la figura: los indicadores de corto plazo se subordinan a la custodia de largo plazo}
A la izquierda, los indicadores de corto plazo incluyen participación de mercado, margen de utilidad, liderazgo tecnológico y volumen de ventas; a la derecha, los objetivos del custodio incluyen la marca, la experiencia del cliente, la solidez financiera y la capacidad de innovación. El CEO de Mercedes-Benz no puede optimizar solo un trimestre ni perseguir únicamente una etiqueta tecnológica.
\end{knowledgebox}

\subsection{No quedar atado a ninguna tecnología}

Lo anterior trató los indicadores del custodio de la marca; esta sección llega a una capa más profunda de la identidad de la marca. Al final de la entrevista, Zhang Xiaojun preguntó si los automóviles de gasolina desaparecerán por completo y, si ese día llega, qué será Mercedes-Benz. La respuesta de Kang Linsong fue que Mercedes-Benz avanza hacia las cero emisiones, pero que no quedará atada a ninguna tecnología específica. La tecnología siempre cambia, y Mercedes-Benz siempre ha permanecido. Esta respuesta se parece mucho al principio de fondo de una marca centenaria: no amarrar la identidad al motor, a la batería, a la pantalla ni a alguna función de software, sino vincularla a promesas más estables.

\begin{importantbox}{La tecnología cambia; la promesa de la marca no puede cambiar}
Para Mercedes-Benz, el motor de combustión, la plataforma eléctrica, la cabina con IA y la conducción autónoma son tecnologías de una etapa. Lo que la marca realmente debe preservar es la seguridad, la calidad, la ingeniería, la experiencia, la estética y la confianza del cliente. El objetivo de la transformación no es preservar la tecnología antigua, sino hacer que estas promesas atraviesen el nuevo ciclo tecnológico.
\end{importantbox}

\subsection{Lista de riesgos de la transformación}

Esta sección condensa la perspectiva del CEO en una lista de riesgos. Cuando una empresa centenaria fracasa en su transformación, no suele ser porque no vea las tendencias, sino porque varias restricciones tiran en direcciones opuestas: por temor a dañar a la base de clientes existente, transforma con demasiada lentitud; por temor a perder nuevas tecnologías, dispersa demasiado la inversión; por temor a la complejidad organizacional, decide con demasiada lentitud; por temor a la presión financiera, reduce la innovación. La respuesta de Kang Linsong puede entenderse como un equilibrio dinámico entre estos riesgos.

\noindent\begin{tabularx}{\textwidth}{p{0.20\textwidth}p{0.34\textwidth}X}
\toprule
Riesgo & Manifestación & Forma de gobierno correspondiente \\
\midrule
Transformación demasiado lenta & Aferrarse solo a la combustión y al lujo tradicional & Invertir en MBOS, IA, plataformas eléctricas e investigación y desarrollo en China. \\
Transformación demasiado brusca & Daño al valor de reventa, a la marca y a las utilidades & Conservar la variedad de productos y proteger el valor para los clientes de gama alta. \\
Pérdida de control tecnológico & Fragmentación tecnológica entre socios externos & Usar MBOS para integrar la arquitectura y controlar la experiencia. \\
Lentitud organizacional & Demasiados niveles jerárquicos en una gran empresa & Simplificar los procesos y elevar la eficiencia y el espíritu emprendedor. \\
Secuestro por indicadores de corto plazo & Perseguir solo la participación de mercado o las utilidades trimestrales & Buscar el equilibrio de largo plazo desde el papel de custodio de la marca. \\
\bottomrule
\end{tabularx}

\subsection{Resumen de la sección}

La tarea del CEO en la transición es convertir la incertidumbre en una capacidad organizacional de corrección continua. Mercedes-Benz debe recuperar el espíritu pionero de su fundación y, al mismo tiempo, mantener la promesa de largo plazo de una marca centenaria.
